% GENERATED FILE. Edit canonical lessons, then run npm run generate:books.
% canonical-source-hash: fnv1a64:3a1a2024b7631300
% canonical-lessons: ZH-C245-tiwen, ZH-C245-chuanran, ZH-C245-kouzhao, ZH-C245-xinqing, ZH-C245-baozhong

\chapter{Take Care}
\label{ch:zh-take-care}
\hlchaptermodality{245}

\begin{chapteropening}
\textbf{By the end of this chapter:} \emph{I can explain why I take my temperature, keep a mask and stay home, and tell someone to look after themselves.}

Complete the last lesson of chapter 245: I can explain why I take my temperature, keep a mask and stay home, and tell someone to look after themselves.
\end{chapteropening}

\section[tǐwēn]{tǐwēn — body temperature}

\label{lesson:ZH-C245-tiwen}



\begin{quote}
Before the new one: say the Mandarin for in case, then the Mandarin for to carry on.
\end{quote}

\subsection*{You'll want to know: tǐwēn}
\emph{tǐwēn} — \textquotedblleft{}body temperature\textquotedblright{}.

\begin{quote}
\emph{nǐ de tǐwēn yǒudiǎnr gāo} — Your temperature's a bit high.
\end{quote}

Its characters come later; for now it is a word for the ear.

\begin{cousinweb}
\emph{tǐwēn} is built from \zh{体} 'body' and \emph{wēn} 'warm', both inherited from Middle Chinese, \emph{thejX} and \emph{'won} (Guangyun, in Baxter's transcription).
\end{cousinweb}

\begin{grammarlens}[title={What you've built}]
One more word for giving reasons.
\end{grammarlens}

\subsection*{Your turn}
\begin{itemize}
\raggedright
  \item \emph{Say it:} \emph{tǐwēn}
  \item \emph{Say it:} \emph{tǐwēn}, once more
  \item \emph{Say it:} \emph{nǐ de tǐwēn yǒudiǎnr gāo} — Your temperature's a bit high.
  \item \emph{Recall:} say the Mandarin for in case, then the Mandarin for to carry on, then say \emph{tǐwēn} again
\end{itemize}

\begin{tcolorbox}[breakable,colback=teal!4,colframe=teal!35!black,title={Before you move on}]
What does tǐwēn mean? (\textquotedblleft{}Body temperature\textquotedblright{}.) Say it once more.
\end{tcolorbox}
\section[chuánrǎn]{chuánrǎn — to pass on (an illness); to be catching}

\label{lesson:ZH-C245-chuanran}



\begin{quote}
Before the new one: say the Mandarin for to carry on, then the Mandarin for body temperature.
\end{quote}

\subsection*{You'll want to know: chuánrǎn}
\emph{chuánrǎn} — \textquotedblleft{}to pass on (an illness); to be catching\textquotedblright{}.

\begin{quote}
\emph{yīnwèi gǎnmào huì chuánrǎn suǒyǐ wǒ bú qù shàngbān} — Colds are catching, so I'm not going in to work.
\end{quote}

Its characters come later; for now it is a word for the ear.

\begin{cousinweb}
\emph{chuánrǎn} is built from \emph{chuán} 'to pass on' and \emph{rǎn} 'to dye', both inherited from Middle Chinese, \emph{drjwen} and \emph{nyemX} (Guangyun, in Baxter's transcription), which glosses the first 'to turn' and the second 'to dye'.
\end{cousinweb}

\begin{grammarlens}[title={What you've built}]
One more word for giving reasons.
\end{grammarlens}

\subsection*{Your turn}
\begin{itemize}
\raggedright
  \item \emph{Say it:} \emph{chuánrǎn}
  \item \emph{Say it:} \emph{chuánrǎn}, once more
  \item \emph{Say it:} \emph{yīnwèi gǎnmào huì chuánrǎn suǒyǐ wǒ bú qù shàngbān} — Colds are catching, so I'm not going in to work.
  \item \emph{Recall:} say the Mandarin for to carry on, then the Mandarin for body temperature, then say \emph{chuánrǎn} again
\end{itemize}

\begin{tcolorbox}[breakable,colback=teal!4,colframe=teal!35!black,title={Before you move on}]
What does chuánrǎn mean? (\textquotedblleft{}To pass on (an illness); to be catching\textquotedblright{}.) Say it once more.
\end{tcolorbox}
\section[kǒuzhào]{kǒuzhào — face mask}

\label{lesson:ZH-C245-kouzhao}



\begin{quote}
Before the new one: say the Mandarin for body temperature, then the Mandarin for to be catching.
\end{quote}

\subsection*{You'll want to know: kǒuzhào}
\emph{kǒuzhào} — \textquotedblleft{}face mask\textquotedblright{}.

\begin{quote}
\emph{nǐ yǒu kǒuzhào ma} — Have you got a face mask?
\end{quote}

Its characters come later; for now it is a word for the ear.

\begin{cousinweb}
\emph{kǒuzhào} is built from \zh{口} 'mouth' and \emph{zhào} 'a cover', both inherited from Middle Chinese, \emph{khuwX} and \emph{traewH} (Guangyun, in Baxter's transcription), which glosses the second as a bamboo basket for catching fish, so how it came to mean a cover is not settled.
\end{cousinweb}

\begin{grammarlens}[title={What you've built}]
One more word for giving reasons.
\end{grammarlens}

\subsection*{Your turn}
\begin{itemize}
\raggedright
  \item \emph{Say it:} \emph{kǒuzhào}
  \item \emph{Say it:} \emph{kǒuzhào}, once more
  \item \emph{Say it:} \emph{nǐ yǒu kǒuzhào ma} — Have you got a face mask?
  \item \emph{Recall:} say the Mandarin for body temperature, then the Mandarin for to be catching, then say \emph{kǒuzhào} again
\end{itemize}

\begin{tcolorbox}[breakable,colback=teal!4,colframe=teal!35!black,title={Before you move on}]
What does kǒuzhào mean? (\textquotedblleft{}Face mask\textquotedblright{}.) Say it once more.
\end{tcolorbox}
\section[xīnqíng]{xīnqíng — mood; frame of mind}

\label{lesson:ZH-C245-xinqing}



\begin{quote}
Before the new one: say the Mandarin for to be catching, then the Mandarin for a face mask.
\end{quote}

\subsection*{You'll want to know: xīnqíng}
\emph{xīnqíng} — \textquotedblleft{}mood; frame of mind\textquotedblright{}.

\begin{quote}
\emph{shēngbìng le xīnqíng bù hǎo} — Being ill puts me in a bad mood.
\end{quote}

Its characters come later; for now it is a word for the ear.

\begin{cousinweb}
\emph{xīnqíng} is built from \zh{心} 'heart' and \emph{qíng} 'feeling', both inherited from Middle Chinese, \emph{sim} and \emph{dzjeng} (Guangyun, in Baxter's transcription), which explains the second, after an older book, as what a person desires.
\end{cousinweb}

\begin{grammarlens}[title={What you've built}]
One more word for giving reasons.
\end{grammarlens}

\subsection*{Your turn}
\begin{itemize}
\raggedright
  \item \emph{Say it:} \emph{xīnqíng}
  \item \emph{Say it:} \emph{xīnqíng}, once more
  \item \emph{Say it:} \emph{shēngbìng le xīnqíng bù hǎo} — Being ill puts me in a bad mood.
  \item \emph{Recall:} say the Mandarin for to be catching, then the Mandarin for a face mask, then say \emph{xīnqíng} again
\end{itemize}

\begin{tcolorbox}[breakable,colback=teal!4,colframe=teal!35!black,title={Before you move on}]
What does xīnqíng mean? (\textquotedblleft{}Mood; frame of mind\textquotedblright{}.) Say it once more.
\end{tcolorbox}
\section[bǎozhòng]{bǎozhòng — to look after yourself; to take care of oneself}

\label{lesson:ZH-C245-baozhong}



\begin{quote}
Before the new one: say the Mandarin for a face mask, then the Mandarin for mood.
\end{quote}

\subsection*{You'll want to know: bǎozhòng}
\emph{bǎozhòng} — \textquotedblleft{}to look after yourself; to take care of oneself\textquotedblright{}.

\begin{quote}
\emph{nǐ shēngbìng le yào bǎozhòng} — You're ill, so look after yourself.
\end{quote}

Its characters come later; for now it is a word for the ear.

\begin{cousinweb}
\emph{bǎozhòng} is built from \emph{bǎo} 'to keep' and \zh{重} 'heavy', both inherited from Middle Chinese, \emph{pawX} and \emph{drjowngX} (Guangyun, in Baxter's transcription), which glosses the first 'to be responsible for' and 'to keep safe'.
\end{cousinweb}

\begin{grammarlens}[title={What you've built}]
One more word for giving reasons.
\end{grammarlens}

\subsection*{Your turn}
\begin{itemize}
\raggedright
  \item \emph{Say it:} \emph{bǎozhòng}
  \item \emph{Say it:} \emph{bǎozhòng}, once more
  \item \emph{Say it:} \emph{nǐ shēngbìng le yào bǎozhòng} — You're ill, so look after yourself.
  \item \emph{Recall:} say the Mandarin for a face mask, then the Mandarin for mood, then say \emph{bǎozhòng} again
\end{itemize}

\begin{tcolorbox}[breakable,colback=teal!4,colframe=teal!35!black,title={Before you move on}]
What does bǎozhòng mean? (\textquotedblleft{}To look after yourself; to take care of oneself\textquotedblright{}.) Say it once more.
\end{tcolorbox}
